\unit{Atomic clocks and precision spectroscopy}
\unitref{Week 13; 2025 Q1; atomic applications}
\subsection{What makes an atom a clock?}
A clock is an oscillator whose phase
can be compared against a reference.
An atomic transition provides a frequency
associated with an exceptionally stable
energy difference:
\[
\boxed{\nu_0=\frac{E_{\rm upper}-E_{\rm lower}}{h}.}
\]
A laser or microwave oscillator is steered
to keep it on resonance with this
frequency. The atom does not generate
the counted electrical clock pulses
by itself: an interrogation sequence
and feedback system are needed.

\subsection{The caesium definition of the SI second}
The unperturbed ground-state hyperfine
transition of \(^{133}\mathrm{Cs}\) defines
the second, with frequency
\[
\boxed{\nu_{\rm Cs}=9\,192\,631\,770\ \mathrm{Hz}.}
\]
Caesium's valence ground term is
\(6s\,{}^2S_{1/2}\), so \(J=1/2\),
while the nuclear spin is \(I=7/2\).
Hyperfine coupling gives \(F=3,4\).
The clock transition is the
\[
\boxed{\ket{F=3,m_F=0}
\longleftrightarrow\ket{F=4,m_F=0}}
\]
transition. The choice \(m_F=0\) makes
the first-order Zeeman shift vanish
near zero field, reducing a major
source of frequency error.

\begin{worked}{Compute the caesium hyperfine interval}
With \(I=7/2,J=1/2\),
\(I(I+1)=63/4\) and \(J(J+1)=3/4\).
The magnetic-dipole hyperfine energy is
\[
E_F=\frac A2
[F(F+1)-I(I+1)-J(J+1)].
\]
Substitute \(F=3,4\):
\[
E_3=\frac A2(12-66/4)
=-\frac{9A}{4},\qquad
E_4=\frac A2(20-66/4)
=\frac{7A}{4}.
\]
Therefore \(E_4-E_3=4A\), so
\[
\boxed{A/h=\nu_{\rm Cs}/4
=2.2981579425\ \mathrm{GHz}.}
\]
The degeneracy check is
\(7+9=16=(2I+1)(2J+1)\).
\end{worked}

\subsection{Ramsey spectroscopy: letting phase accumulate}
In a Ramsey experiment a first
\(\pi/2\) pulse prepares an equal
superposition of clock states.
The atomic components then acquire
a relative phase during free evolution
for time \(T\). A second \(\pi/2\)
pulse converts that phase into
population. In the idealised
two-level model, the excited
population oscillates with detuning:
\[
P_e\sim\tfrac12[1+\cos(\Delta T+\phi_0)].
\]
The fringe width scales approximately
as \(1/T\). Long coherent interrogation
therefore permits a narrow frequency
discriminator, provided the atom
and the oscillator stay coherent.

\subsection{Why optical clocks can do better}
The quality factor of a spectral
reference is
\[
\boxed{Q=\frac{\nu_0}{\Delta\nu}}.
\]
A high carrier frequency and narrow
linewidth can yield a very high \(Q\).
For example, strontium optical clocks
use a transition near \(429\) THz,
vastly higher than the \(9.19\) GHz
caesium hyperfine carrier.
Optical clocks usually use specially
chosen metastable, narrow transitions
rather than bright, rapidly
decaying laser-cooling lines.

\subsection{Trapping without shifting the clock}
An optical lattice localises atoms,
reducing motional Doppler effects.
However, trapping light also shifts
atomic energy levels through the
AC Stark effect. At a
\emph{magic wavelength}, the
differential light shift of the
two clock states is approximately
zero, permitting confinement
without substantially shifting
the measured transition frequency.

\subsection{Stability and accuracy are different}
Stability describes the scatter
of clock frequency measurements
over an averaging time; accuracy
describes how closely the corrected
frequency matches the unperturbed
atomic reference. Quantum projection
noise, technical oscillator noise
and finite atom number affect
stability. Zeeman/Stark shifts,
black-body radiation, collisions
and gravitational redshift are
among the corrections important
for accuracy. A clock can be
very stable yet systematically wrong.

\subsection{Why clocks are physics experiments}
Comparing clocks based on different
species or transitions provides tests
of relativity and searches for
changes in fundamental constants.
Relativistic gravitational redshift
means clocks at different gravitational
potentials tick at measurably different
rates. Precise atomic spectroscopy
therefore connects quantum
structure to timekeeping, geodesy
and fundamental physics.

\begin{checkpoint}
A clock should have a long-lived,
narrow transition and be insensitive
to external perturbations.
The atomic transition, interrogation
sequence, local oscillator, feedback
and systematic corrections all
determine clock performance.
\end{checkpoint}
